\begin{abstract}
Cross-lingual claim verification systems frequently convert non-English evidence texts into intermediate English representations prior to automated veracity classification. However, translating evidence introduces a critical tension between exploiting high-capacity English language models and introducing semantic distortion. In this study, we present a controlled empirical investigation of verifier prediction sensitivity under Hindi-to-English translation shift on authentic claim-evidence pairs from the X-Fact benchmark ($N=266$: 103 development, 163 held-out test). Holding the claim, pre-retrieved evidence snippets, and cross-lingual verifier (\texttt{mDeBERTa-v3-base}) fixed, we evaluate veracity decisions across four paired conditions: \emph{No-Check} (machine-translated English evidence), \emph{Direct Multilingual} (original Hindi evidence), \emph{Always-Check} (unconditional posterior consensus), and a risk-calibrated \emph{Selective-Check} cascade tuned on development data ($\tau^*=0.650$).

We observe an overall paired prediction disagreement rate of 63.80\% (104/163; 95\% CI: [56.19\%, 70.78\%]) between original and translated evidence representations. When translated evidence predicts \textsc{Supported}, evaluating the original Hindi evidence flips the prediction to \textsc{Contradicted} in 48.44\% of cases (31/64). Direct Multilingual verification achieves 42.94\% accuracy (70/163), significantly outperforming the No-Check baseline (28.22\%, 46/163; McNemar exact $p = 0.008772$, $b=27, c=51$). However, risk-calibrated Selective-Check does not surpass No-Check accuracy (28.22\%, 46/163; coverage 7.98\%), and paired permutation tests on class-balanced Macro-F1 reveal no statistically significant difference across conditions ($p = 0.1898$ and $p = 0.1180$), reflecting the high variance of harmonic metrics under severe natural class skew (87.12\% \textsc{Contradicted}, 0.0\% gold \textsc{Supported}). We transparently document the exploratory nature of this evaluation, the prior inspection of the test split, and the necessity of pending bilingual human validation before causal translation claims can be substantiated.
\end{abstract}
